% Original PDF 29 / printed 19. All of Box 2.1, not included in accepted sample.
% Source box is represented using the existing user's template color, not a new style.
% block: homological-algebra-summary-heading page:29 type:ordinary-paragraph
\begin{tcolorbox}[enhanced,breakable,colback=KuangNei,frame hidden,boxrule=0pt,sharpish corners,before skip=10pt,after skip=10pt]
\refstepcounter{booksummary}\label{box:homological-algebra-summary}
\noindent\textbf{Box~\thebooksummary:} \BT{Summary of Chapter~\ref{chap:homological-algebra} on homological algebra.}{Chapter~\ref{chap:homological-algebra} 同调代数内容摘要.}\par

% block: summary-chain-complex-intro page:29 type:ordinary-paragraph
\BT{A \emph{chain complex} is a sequence of vector spaces and operators}{\emph{链复形}是如下一列向量空间与算子:}

% block: summary-chain-sequence page:29 type:equation
\[
\cdots\to V_{k+1}\xrightarrow{\partial_{k+1}}V_k\xrightarrow{\partial_k}V_{k-1}\to\cdots
\]

% block: summary-chain-spaces page:29 type:ordinary-paragraph
\BT{with successive operators composing to zero. It gives rise to three vector spaces at each level $k$:}{其中相邻算子的复合为零. 它在每个层次 $k$ 上产生三个向量空间:}
\begin{itemize}
\item \BT{$\mathcal Z_k$, the null space of the boundary operator $\partial_k$, consisting of the \emph{cycles};}{$\mathcal Z_k$, 即边界算子 $\partial_k$ 的零空间, 由\emph{闭链}组成;}
\item \BT{$\mathcal B_k$, the range of $\partial_{k+1}$, the \emph{boundaries};}{$\mathcal B_k$, 即 $\partial_{k+1}$ 的像空间, 由\emph{边界链}组成;}
\item \BT{$\mathcal H_k=\mathcal Z_k/\mathcal B_k$ the \emph{homology space}.}{$\mathcal H_k=\mathcal Z_k/\mathcal B_k$, 即\emph{同调空间}.}
\end{itemize}

% block: summary-simplicial-example page:29 type:ordinary-paragraph
\BT{An important example is the \emph{simplicial chain complex} associated to a triangulation or more general simplicial complex. The dimensions of its homology spaces are the Betti numbers of the underlying space.}{一个重要例子是与剖分或更一般的单纯复形相伴的\emph{单纯链复形}. 它的同调空间的维数是底空间的 Betti 数.}

% block: summary-cochain-intro page:29 type:ordinary-paragraph
\BT{A \emph{cochain complex} has the indexing reversed,}{\emph{上链复形}的指标次序相反,}

% block: summary-cochain-sequence page:29 type:equation
\[
\cdots\to V^{k-1}\xrightarrow{d^{k-1}}V^k\xrightarrow{d^k}V^{k+1}\to\cdots,
\]

% block: summary-cochain-objects page:29 type:ordinary-paragraph
\BT{giving rise to \emph{cocyles}, \emph{coboundaries}, and \emph{cohomology}.}{由此产生\emph{闭上链}, \emph{上边界链}与\emph{上同调}.}

% block: summary-duality page:29 type:ordinary-paragraph
\BT{The dual spaces and dual maps of a chain complex form a cochain complex.}{链复形的对偶空间和对偶映射构成一个上链复形.}

% block: summary-de-rham-intro page:29 type:ordinary-paragraph
\BT{Another cochain complex is the \emph{de Rham complex} which, on a three-dimensional domain, can be written}{另一个上链复形是 \emph{de Rham 复形}, 在三维区域上, 它可以写成}

% block: summary-de-rham-sequence page:29 type:equation
\[
0\to C^\infty(\Omega)\xrightarrow{\operatorname{grad}}C^\infty(\Omega;\mathbb R^3)\xrightarrow{\operatorname{curl}}C^\infty(\Omega;\mathbb R^3)\xrightarrow{\operatorname{div}}C^\infty(\Omega)\to0.
\]

% block: summary-de-rham-theorem page:29 type:ordinary-paragraph
\BT{The de Rham map is a \emph{cochain map} from the de Rham complex to the dual of the simplicial chain complex and, as such, induces a map on cohomology. De Rham's theorem states that the induced map is an isomorphism.}{de Rham 映射是从 de Rham 复形到单纯链复形的对偶复形的一个\emph{上链映射}, 因而在上同调上诱导一个映射. De Rham 定理断言, 这一诱导映射是同构.}
\end{tcolorbox}
